\section{Introduction}\label{sec:introduction}

The pilot of a remotely operated vehicle (ROV) works through a
camera. Water takes
the red from the picture, scatters light into the line of sight and
lowers contrast
with distance \cite{jaffe1990,schettini2010}, and the vehicle's own lights are
scattered back at the camera \cite{abdullah2026}; whatever is done to
the picture to
compensate must be done between the camera and the pilot without
holding it back,
because delay in the picture degrades visually guided work
\cite{mackenzie1993,nankaku2022,chen2007}. Small inspection-class
vehicles carry a
single-board computer, an autopilot and an Ethernet tether
\cite{vonbenzon2022}, and
that computer is the place where such processing has to fit.

Three bodies of work bear on this, and Section~\ref{sec:related}
reviews them. Image
enhancement for underwater pictures is a large field whose recent
methods are learned
and run in real time on graphics processors at low resolution; the model-free
operators that are cheap enough for a small processor can flicker when their
parameters are estimated anew for every frame, and the remedies for
flicker are not real-time.
Video pipelines built from standard parts report delays of 60 to
200~ms from encoder
to decoder or from scene to display, of which the instrumented stages
explain a few
milliseconds; for a small ROV observed through a network camera, 850~ms has been
measured \cite{kaknjo2018}. And vehicles that add perception pass frames between
processes through robotics middleware, whose cost for messages of the size of an
image has been measured on single-board computers and is high, or
send the video to
the surface and process it there.

This paper describes electric-eye, open-source software in two programs for the
video of an ROV between its camera and its pilot, and makes the following
contributions.
\begin{itemize}
  \item A continuous visibility correction for the pilot's picture
    (cast removal,
    contrast stretch and local histogram equalization) whose
    parameters, rather than
    its pixels, are smoothed over time, at 0.55~ms per
    1280$\times$720 frame on a
    desktop processor, beside a chain of image operators with NEON and
    multi-threaded paths that are required to reproduce the scalar
    output byte for
    byte (Sections~\ref{sec:operators} and~\ref{sec:autovis}).
  \item A frame bus in shared memory with a documented layout, through which any
    program on the machine reads the camera's frame and the pilot's picture with
    their capture times, and whose writer never waits for a reader
    (Section~\ref{sec:bus}).
  \item A transport of the video across the tether as H.264 over RTP
    in which two
    changes at the receiver, neither of them to the encoder or the
    network, reduced
    the latency from bus to bus over loopback from 70.4 to 3.8~ms at
    640$\times$480
    and from 74.8 to 8.3~ms at 1280$\times$720; and the option to
    send a camera's
    own H.264 with no encoder on the vehicle (Section~\ref{sec:transport}).
  \item The freedom to run the image processing on the vehicle or at the surface
    with the same programs (Section~\ref{sec:overview}).
\end{itemize}
The measurements are from the bench. The software has not yet been
run on a vehicle,
the visibility correction has been evaluated on synthetic frames
only, and no pilot
has been studied; Section~\ref{sec:discussion} returns to what
follows from that.

Section~\ref{sec:related} reviews related work,
Section~\ref{sec:methods} describes
the design and the methods of verification and measurement,
Section~\ref{sec:results} gives the results,
Section~\ref{sec:discussion} discusses
them, and Section~\ref{sec:conclusion} concludes.
